\begin{abstract}
\noindent
The Zenodo record 10.5281/zenodo.22883643, \emph{SU(2) lattice Yang--Mills: vacuum, energy, spectral bounds and continuum investigations}, is the Yang--Mills part of the owner's research workbench (\texttt{KokunoYumeto/yang-mills-interacting-workbench}). It holds 43 files from five editions, dated 8 to 21 September 2026. The record's front file, \texttt{README.md}, is still the one written for the first edition of 8 September.

This reader says:
\begin{itemize}
\item which of the 43 files is current and which are superseded;
\item what each document claims, and on which regime;
\item what is proved, with proofs;
\item how far each statement was checked.
\end{itemize}

The central result (Section~\ref{sec:gap}) concerns the Kogut--Susskind Hamiltonian of $SU(2)$ lattice gauge theory on open cubic boxes. It is an explicit lower bound on the spectral gap above the ground state, on the whole gauge-invariant space, uniform in the box size and the lattice spacing, for bare coupling $g^2\ge3.6836$ in the workbench's normalisation. The threshold rests on computed coefficient bounds of orders 3 to 5. With the order-2 inputs recomputed in the first-pass audit, the same argument gives a gap of at least $1.52\kappa$ for $g^2\ge3.9781$, which includes the workbench's benchmark $g^2=4$.

Three consequences of the record's material are added here:
\begin{itemize}
\item the gap persists at complex coupling, so the ground state depends analytically on the coupling on a disk whose radius does not depend on the box;
\item an upper bound on the gap: for $g^2\ge13$, for example, the gap lies between $2.90\kappa$ and $3.0003\kappa$;
\item a gap bound on the full space of states, not only the gauge-invariant one.
\end{itemize}

Section~\ref{sec:heat} states the heat-flow results of the latest edition:
\begin{itemize}
\item every constant of the box-independent heat remainder is re-derived here, and its prefactor $67896/169$ is improved to $2829/13$ by a sharper estimate of the mean;
\item the benchmark table for relaxation into the whole complementary physical space is recomputed exactly from the stated inputs, and then sharpened.
\end{itemize}
Sections~\ref{sec:earlier}--\ref{sec:side} map the earlier editions and the Jacobian-map, Navier--Stokes and $S^6$ material, and Section~\ref{sec:negative} states the obstructions with their exact scope.

All results hold at fixed lattice spacing and strong bare coupling. The workbench does not claim a continuum mass gap, and nothing here bears on that problem.
\end{abstract}

\section*{About this reader}\phantomsection\label{sec:about}
\addcontentsline{toc}{section}{About this reader}

\paragraph{The workbench, and who made it.}
\begin{itemize}
\item The record's \texttt{README.md} (``Source history, corrections and attribution'') says that the human participant proposed the research directions and possible connections, and that Codex developed, checked and corrected the written constructions. It attributes the public stewardship to KokunoYumeto, the owner of the workbench.
\item The repository's README describes the workbench as open and AI-assisted, including work with ChatGPT~5.6 Sol and GPT-6 Astra.
\item The workbench credits:
\begin{itemize}
\item the exponential form of the vacuum, and its character and Casimir framework, to Schütte, Zheng and Hamer \cite{SchutteZhengHamer};
\item the Fourier-algebra estimates to Eymard \cite{Eymard};
\item semigroup generation to Lumer and Phillips \cite{LumerPhillips};
\item the creation method behind the quantum line's first volume-uniform gap to Yarotsky \cite{Yarotsky}, with Baez's spin-network framework \cite{Baez} as its reference for the physical space (quantum line, p.~19; the record's \texttt{README.md} lists both).
\end{itemize}
\item Two inputs of side branches are credited to Levent Alpöge by the workbench's attribution file (\texttt{ATTRIBUTION.md}; record file \texttt{ALPOGE\_ATTRIBUTION\_20260921.md}). One is a polynomial map that the file identifies with the counterexample to the Jacobian conjecture announced by Alpöge on 20 July 2026, whose announcement, as the file records, credits Akhil for the question and Fable for the work. The other is the construction of a complex threefold, circulated by Alpöge and, as the file says, produced with Claude \cite{AlpogeS6}. The file says that neither is used as a hypothesis of the 21 September heat and volume proofs. That the fifth-reference results (F1--F42, H1--H34) do not use them either follows from reading those proofs (Sections~\ref{sec:gap} and~\ref{sec:heat}).
\item The Navier--Stokes material rests on OpenAI's manuscript \emph{Finite Time Blowup for Navier--Stokes} \cite{OpenAINS}, which the workbench cites and does not claim to have verified.
\end{itemize}

\paragraph{What was read.}
\begin{itemize}
\item The repository at commit \texttt{fa79faf} (21 September 2026), the head of \texttt{origin/main} on 26 September.
\item The record's 43 files. Twenty-three of them are byte-identical to files of that commit (checked by MD5). The other twenty were downloaded from the record on 26 September, except the 87~MB Overleaf archive (file \texttt{08\_}), and their MD5 checksums equal the record's.
\item The two proof manuscripts of the latest edition in full: \texttt{FIFTH\_REFERENCE.md} (F1--F42) and \texttt{HEAT\_AND\_COMPLEMENT.md} (H1--H34).
\item For the other documents, their scope and conclusion sections, and the passages cited below.
\end{itemize}
The audit notes \note{21\_}, \note{28\_}, \note{34\_}, \note{46\_}, \note{46a\_} and \note{47\_}, and the reader of four workbenches \cite{WorkbenchReader}, are on the branch \texttt{claude/claude-ab-grind-20260925} of \texttt{KokunoYumeto/zeta-function-research-reader}, in \texttt{contrib/claude-ab/grind\_20260925/}.

\paragraph{How it was checked.}
\begin{itemize}
\item Section~\ref{sec:gap} comes from the Yang--Mills section of the reader of four workbenches, which was refereed three times, after a referee pass on its notes (Appendix~\ref{app:verification}).
\item The new material of Section~\ref{sec:heat} was checked by this reader's author, with a new script. It recomputes the constants from the manuscript's formulas and the stated inputs. The constants $8$, $32$, $\frac18$, $30$ and $48$ are derived in the text.
\item This version includes the findings of one independent referee pass, by a Claude instance that wrote none of the text. It read the reader for overstatement, for understatement and missed generality, and for missed or implied results. Each finding was verified before it was applied, by reading the record at the cited place or by an independent script (Appendix~\ref{app:verification}).
\item Everything else carries a status label.
\end{itemize}

\paragraph{Status labels.}
\begin{itemize}
\item \textsf{Proved here}: the proof is given in full in this reader.
\item \textsf{Generalised here}: the record's statement is proved here under weaker hypotheses, or with a sharper constant, than the record states.
\item \textsf{Audited}: the workbench's proof was read in full, every step was checked, and independent checks were run.
\item \textsf{Audited in part}: the parts named were checked; the rest is located.
\item \textsf{Reproduced}: the finite content of a statement (its constants, tables or counts) was recomputed by an independent script from the stated inputs; the general argument was not re-derived.
\item \textsf{Conditional}: proved, given inputs that were not checked here; the inputs are named.
\item \textsf{Located}: the statement and its proof were found and are described; nothing was checked.
\end{itemize}
The workbench itself separates written analytic arguments from finite exact-arithmetic certificates and from numerical diagnostics, and says so at each result; its own qualifications are quoted where they matter.
